\documentclass[12pt,a4paper]{article}
\usepackage[margin=28mm]{geometry}
\usepackage[T1]{fontenc}
\usepackage{lmodern,amsmath,amssymb,booktabs,graphicx,microtype,xurl,setspace}
\usepackage[numbers,sort&compress]{natbib}
\usepackage[hidelinks]{hyperref}
\usepackage{xr}
\externaldocument{Supplement_labels}
\graphicspath{{../Figures/}}
\newcommand{\R}{\mathbb R}
\input{experimental_sections}
\doublespacing
\begin{document}
\begin{center}
{\Large\bfseries Global loss bounds and one-bit recovery for direction-normalized half-moons\par}
\vspace{1em}
{\large Meshaal Mouawad\par}
Department of Electrical and Computer Engineering, Mississippi State University\\
Mississippi State, MS 39762, USA\\
Corresponding author: \texttt{mm4922@msstate.edu}
\end{center}
\vspace{1em}
\noindent\textbf{Abstract.}
Direction normalization retains the orientation of an observation from a shared origin while discarding its distance. If that distance determines the optimal class decision, a classifier trained on the normalized representation cannot recover the lost distinction. We study this mechanism in a continuous half-moon family and determine whether moving the normalization origin can prevent the resulting loss. An exact ray-cancellation identity expresses additional Bayes error through class-preference mass merged by direction. A uniform distant-origin limit to orthogonal-projection risk and geometric criteria establish positive loss over every origin for specified disjoint and overlapping half-moon configurations. Constructive certificates give explicit whole-plane lower bounds. For one disjoint midpoint geometry, adding a common-threshold near/far radius bit restores every class decision, with a sharp threshold interval and measurement-error tolerance. Frozen neural studies contain 500 candidate fits and 250 selected evaluations. In the independent repair experiment, direction-only networks made 54,362 errors among 500,000 observations; the one-bit networks made three. The analysis separates representation loss from learning error and establishes a minimal radial decision repair in the proved geometry.

\medskip\noindent\textbf{Keywords:}
Direction normalization; Bayes risk; spatial sign; half-moon classification; information preservation; radial quantization; one-bit decoder.

\section{Introduction}
A preprocessing transformation changes what a classifier can learn before training begins. Many transformations intentionally suppress aspects of an observation, such as measurement scale, that may be irrelevant to a particular task. When an attribute treated as nuisance variation also carries information needed for classification, its removal creates an error that no subsequently trained classifier can avoid. Distinguishing this information-induced error from imperfect model fitting is central to evaluating representations rather than merely comparing the performance of learning algorithms.

Direction normalization provides a concrete instance. Given an observation and a common reference origin, it retains the unit vector pointing from the origin to the observation and discards the distance along that direction. Such a representation is meaningful when orientation, rather than magnitude, should govern the task; related spatial-sign and normalized-feature transformations have established statistical and classification uses \citep{Serneels2006,Graf2003}. Yet two observations on the same ray receive identical features. If those observations require different decisions, a classifier seeing only direction has no variable with which to separate them. Whether this collapse causes a positive population error, rather than an isolated geometric coincidence, depends on the underlying class distributions.

The half-moon problem makes the distinction explicit. Its two labels correspond to thick semicircular regions: the upper half of a ring and a lower half-ring translated horizontally and vertically. Their relative positions produce both separable and overlapping configurations while keeping the geometry analytically tractable. At vertical displacement \(h=0.5\), choose the common midpoint \(c=(0.5,0.25)\). The class-0 observation \(x_0=(1,0.25)\) and class-1 observation \(x_1=(2,0.25)\) belong to their respective half-annuli, yet normalize to the same direction \((1,0)\). Their distances from the midpoint are \(0.5\) and \(1.5\). This constructed example identifies the discarded distinction; a single near/far radius bit at threshold \(0.8\) distinguishes it. A theorem is needed to show that a nonzero population of decisions is lost, and that changing the origin cannot eliminate the loss.

We quantify the effect using \emph{Bayes risk}, the minimum misclassification probability achievable with an optimal decision rule when the class distributions are known. The risk for original coordinates captures ambiguity intrinsic to the task. The risk for direction-only inputs measures what remains achievable after preprocessing. Their difference is the additional error imposed by the representation, regardless of model architecture or optimizer. A trained network may incur further error because of finite samples, approximation and optimization; that excess must be evaluated against its \emph{own} representation-specific Bayes floor.

Three questions structure this study. First, when does direction normalization remove class information, and can any origin, including one arbitrarily far away, avoid that loss? Second, when direction alone is insufficient, how much radial information is necessary for the class decision, and how robust is the resulting code to measurement error? Third, do neural classifiers fitted under a frozen protocol approach the corresponding representation-specific Bayes risks? Each question requires distinct evidence: mathematical identities and geometric proofs, a constructive decision rule and sharp interval, and independent held-out neural evaluations.

We derive an exact ray-cancellation identity for the loss, prove a uniform limit connecting distant normalization origins to linear projection risk, and construct positive whole-plane lower certificates. For a specified disjoint configuration, one common near/far radius bit selects between two angular separators and correctly labels all points in both continuous class supports. We then compare raw, direction-only, radius-preserving, and one-bit representations in controlled studies comprising 500 candidate fits and 250 selected test evaluations. The combined contribution is a population-level explanation of when preprocessing destroys a decision, why choosing a new origin cannot necessarily repair it, and how a minimal additional measurement can do so in an explicitly proved setting.

\section{Research context and relationship to prior work}
\subsection{Direction and the geometry of normalized features}
The spatial-sign transformation centers an observation and retains its unit direction, an operation studied in robust multivariate analysis \citep{Serneels2006}. Work on normalized feature spaces also addresses how classification procedures interact with angular feature geometry \citep{Graf2003}. Such methods establish meaningful uses of direction representations. They do not by themselves determine whether a given label decision is invariant to radius about every possible origin.

Recent research in \emph{Pattern Recognition Letters} clarifies why geometry and decision rules must be distinguished. Poblete et al.~\citep{Poblete2026} hold generative posterior-mean embeddings fixed and compare radial, Gaussian probabilistic, linear discriminative and tree-based classifiers across eight datasets. Classifier rankings vary with the dataset; their class-centered radial-compatibility descriptor addresses the fit between a decision-rule assumption and representation geometry. This identifies a compatibility question \emph{within} an existing representation. Our question concerns an earlier step: whether a transformation into direction-only features has erased distinctions even an optimal decision rule would require.

Ahmadian and Lindsten~\citep{Ahmadian2025} investigate time-series out-of-distribution detection with radial-kernel scoring and context quantization. Their separate ablations show that a normalized representation performs differently from an unnormalized radial-kernel configuration on the principal activity datasets. The observed difference makes lost magnitude information a plausible task-specific explanation. Their outcome is novelty-detection performance, whereas the present work measures binary Bayes classification risk under known continuous distributions. The studies are complementary in motivation rather than directly comparable numerical benchmarks.

\subsection{Decision sufficiency and representation error}
Silva and Narayanan~\citep{Silva2012} treat signal representation in a Bayes-decision framework and distinguish representational approximation from finite-data estimation. Further work on decision sufficiency and lossless transformations establishes when a compressed representation retains the optimal decision for a specified loss \citep{Silva2024,Gyorfi2023,Silva2025}. The general principle is that information about an observation need not be preserved in full: what matters for classification is whether the optimal label can still be determined.

Direction normalization tests that principle on an explicit family of equivalence classes. All points on the same positive ray from an origin have the same unit direction. Dubois et al.~\citep{Dubois2021} identify unit direction as an invariant to positive scaling when the predictive task satisfies the corresponding invariance condition. If the preferred class changes along some of these rays, the direction map violates the decision-sufficiency condition. The present analysis asks how that violation depends on a moving common origin and supplies an exact probability identity for its cost.

\subsection{Projection boundaries and restricted radial recovery}
Studies of projected classification and spherical models characterize decisions under selected projections or geometric distribution assumptions \citep{Hamsici2007,Hamsici2008,Oh2025}. From increasingly distant origins, rays through a bounded population become nearly parallel, linking radial viewing geometry to orthogonal projections \citep{Orponen2019}. Separability results for planar wedges and strips provide geometric precedents for reasoning about mutually obstructing class regions \citep{Hurtado2001}. These strands motivate the need to handle both finite origins and unbounded sequences of origins before claiming a strictly positive minimum over the entire plane.

A representation can also retain a small task-specific code rather than reconstruct the discarded variable. Bayes-aware quantization and learned threshold research address decision-relevant compression \citep{OehlerGray1995,Senozan2024}. In the present study, the code is deliberately constrained to one common threshold applied to radius and two fixed angular separators. This establishes a precise question: can one bit, used together with direction, recover the correct decision everywhere on the specified supports, and what perturbation of the measured radius can be tolerated?

Taken together, this literature establishes the value of directional representations, the need for decision-aware sufficiency criteria, and mathematical tools for projection and quantization. The specific gap addressed here is an origin-wide population-risk account for nonlinear direction normalization, linked to a minimal and explicitly certified radial decision repair.


\section{Study geometry, representations and inferential design}
\subsection{Why three half-moon arrangements are necessary}
The half-moon family changes the vertical displacement \(h\) of one band while keeping its width and horizontal displacement fixed. The geometry is deliberately continuous: conclusions about population decisions should not depend on the particular finite points sampled for an experiment. The class supports, their probability densities, and the loss incurred by angular projection are defined formally in the next sections.

Three configurations serve different inferential purposes. The \(h=-0.25\) arrangement is a negative control in which suitable origins can preserve all class decisions after direction normalization. At \(h=0.5\), the supports are disjoint, so full coordinates classify them without Bayes error, yet their arrangement forces directional ambiguity from every origin. At \(h=1\), the original regions overlap, creating intrinsic raw-coordinate ambiguity before any normalization; direction adds another component of error. Comparing the three makes it possible to distinguish a representation defect from a genuinely ambiguous source distribution.

\begin{table}[htbp]\centering
\caption{Geometry and intended role of each evaluation configuration. The descriptions summarize proved support properties; they are not inferred from a neural classifier's observed accuracy.}
\label{tab:study-geometry}
\begin{tabular}{p{23mm}p{55mm}p{62mm}}
\toprule
Displacement & Source-space property & Main question\\
\midrule
\(h=-0.25\) & Disjoint; decision-preserving origins exist & When can direction retain the optimal label decision?\\
\(h=0.5\) & Disjoint; no origin eliminates directional loss & How large is preprocessing-induced error when the raw classes are separable?\\
\(h=1\) & Original supports overlap; direction adds loss & How do original ambiguity and preprocessing loss differ?\\
\bottomrule
\end{tabular}
\end{table}

\begin{figure}[htbp]\centering
\includegraphics[width=.96\textwidth]{Fig1_Geometry.pdf}
\caption{The two half-annular class supports under three relative displacements. A shared normalization origin determines which points have the same unit direction; colored rays in the original construction illustrate the geometry of merged decisions. Panel A exhibits an origin preserving the decision, panel B a disjoint arrangement with forced directional overlap, and panel C a configuration with overlapping raw supports. The illustrations motivate, rather than replace, the population-risk proofs.}
\label{fig:geometry}
\end{figure}

\subsection{Sampling laws and what their comparison tests}
The primary experimental distribution samples uniformly with respect to planar area inside each class support, using equal class priors. This means that a region twice as large has twice as much probability within a given class, independent of its radial location. A separate distribution samples the local radius and semicircle angle uniformly; after conversion to planar density, this gives greater density near the inner boundary than the area-uniform law. The two models have identical supports and different probability allocations. Their comparison tests whether the geometric obstruction persists under changes in sampling measure, and how the numerical population risks change.

The frozen neural evaluations use only area-uniform samples. The radius-uniform alternative contributes analytical population-risk results, not an additional neural benchmark. Both are informative: the support geometry establishes qualitative decision-loss conditions whenever class densities are strictly positive on the stated regions, while the exact probability density determines the numerical amount of Bayes error. This distinction is needed to interpret theorems about all origins alongside tables of calculated risks for specified distributions.

\subsection{Representation comparisons and their information content}
The input representation is the independent variable of the neural evaluation. The full-coordinate condition \(R\) retains the two Cartesian measurements. Three direction-only conditions \(D_0,D_M,D_A\) normalize about the origin, the class-center midpoint, and the upper reference point \((0,3)\). Their values differ because the shared center changes which observations lie on the same ray. The direction-plus-radius condition \(P_A\) retains both quantities about the upper origin. Since the original point can be reconstructed from those two quantities and the known center, it has the same population Bayes risk as \(R\), even though a network fitted to the two feature sets need not have identical learned error.

These controls are designed around an information hierarchy. Raw coordinates and exact direction-plus-radius inputs have the full coordinate information. Direction-only features discard the radial component. The one-bit repair condition, evaluated separately, restores a categorical part of the radius without restoring the continuous variable. Comparing population Bayes risks identifies information loss exactly; comparing trained networks shows how those distinct inputs perform under a fixed fitting protocol.

% Main-text fragment; supplementary proofs are in supplement_theory.tex.
% Requires amsmath, amssymb, booktabs and natbib.
\section{Population geometry and the information lost in a directional representation}
\label{sec:origin-theory}

The geometric cases introduced in Fig.~\ref{fig:geometry} are formalized by two continuous half-annular class supports. The displacement $h$ changes the position of the second support; the normalization origin $c$ is a separate choice that determines which observations share a direction. These definitions make it possible to distinguish ambiguity in the original data from the additional Bayes error created by the representation.

\begin{samepage}
Let $a=0.8$, $b=1.2$, $m_0=(0,0)$ and $m_1=(1,h)$, where $h$ is the upward center displacement. The open class supports are
\begin{align}
 S_0&=\{x:a<\|x\|<b,\ x_2>0\},\nonumber\\
 S_1&=\{x:a<\|x-m_1\|<b,\ x_2<h\}.
 \label{eq:halfmoon-supports}
\end{align}
\end{samepage}
Each support has area $H=\pi(b^2-a^2)/2=2\pi/5$. The primary generator draws points uniformly over the area of each band, so its density is $f_i^A=\mathbf1_{S_i}/H$. A second law draws radius and angle uniformly, giving $f_i^R(x)=\mathbf1_{S_i}(x)/[\pi(b-a)\|x-m_i\|]$; it tests whether the geometric conclusions survive a different allocation of probability within the same supports. Neural experiments use equal priors and area-uniform sampling. The geometric positivity results apply to any positive priors $\pi_i$ and densities positive almost everywhere on the respective supports.

The two-point collision in the Introduction demonstrates that normalization is not one-to-one. It does not by itself prove a positive population error, because a point or single ray has zero probability under the continuous sampling laws. The relevant condition is conflict between Bayes decisions across a set of directions with positive angular measure. The following identity measures that conflict from the class densities.

\paragraph{Why one ray can force classification error.}
Direction-only classification must use the same decision for every radius associated with an angle. The class-weighted preference is expressed by the signed density $d=\pi_0f_0-\pi_1f_1$: its positive and negative regions favor classes 0 and 1, respectively. Let $u_\theta=(\cos\theta,\sin\theta)$ and $[z]_+=\max(z,0)$. Integration along the ray from $c$ uses the polar area factor $r$, because an angular slice contains more area at a larger radius. The opposing class-preference masses per unit angle are
\begin{equation}
\begin{aligned}
 A_c(\theta)&=\int_0^\infty[d(c+ru_\theta)]_+r\,dr,\\
 B_c(\theta)&=\int_0^\infty[-d(c+ru_\theta)]_+r\,dr.
\end{aligned}
\label{eq:ray-masses-main}
\end{equation}
Let $R_X^*$ and $R^*_{T_c}$ denote the minimum classification risks from raw coordinates and directions, respectively. Using the masses in~\eqref{eq:ray-masses-main}, the binary Bayes-risk identity yields
\begin{equation}
 \Delta(c):=R^*_{T_c}-R_X^*
 =\int_0^{2\pi}\min\{A_c(\theta),B_c(\theta)\}\,d\theta.
 \label{eq:ray-loss-main}
\end{equation}
The functions $A_c(\theta)$ and $B_c(\theta)$ are signed decision-mass \emph{densities per unit angle}, not probabilities on individual zero-width rays. Their minimum describes opposing preference cancelled at each angle; only integration across angles gives the excess Bayes risk $\Delta(c)$, a probability. Positive loss therefore requires conflicting preferences on a set of directions of positive angular measure. This specializes decision-preservation theory \citep{Silva2012,Silva2024} to a moving normalization origin; Supplementary Section~\ref{supp:bayes} derives the identity.

\paragraph{Theorem 1: loss at every finite origin.}
A collision seen from one origin might disappear when the origin moves. The stronger question is whether some viewpoint separates the two half-rings into different sets of directions. Their geometry answers that question for the positive displacements below: despite disjoint original class regions, every finite viewpoint merges conflicting decisions. For $0<h\leq\sqrt{3/5}$, the supports have disjoint interiors and
\begin{equation}
 R_X^*=0,\qquad \Delta(c)>0
 \quad\text{for every }c\in\mathbb R^2.
 \label{eq:all-origin-disjoint}
\end{equation}
At $h=0$, every origin with $c_2=0$ preserves zero risk. More generally, for $h<0$, origins in the separating strip $h\leq c_2\leq0$ preserve zero risk. At $h=1$, the supports overlap and
\begin{equation}
 R_X^*>0,\qquad \Delta(c)>0
 \quad\text{for every }c\in\mathbb R^2.
 \label{eq:all-origin-overlap}
\end{equation}

\emph{Proof idea.}
Horizontal slicing shows that the disjoint supports satisfy $\sqrt{a^2-h^2}\geq1/5$ in the stated range. Yet each connected half-ring intrudes into the other one's convex hull. Were their directional images disjoint from some origin, a separating angular cone would contain one region's convex hull while excluding the other, contradicting this intrusion. This adapts the geometry of Hurtado et al.~\citep{Hurtado2001}; positive class densities convert the forced angular overlap into positive risk in~\eqref{eq:ray-loss-main}. At $h=1$, the same argument uses connected class-exclusive regions to establish additional loss despite the overlap already present in raw coordinates. Supplementary Sections~\ref{supp:geometry}--\ref{supp:overlap} provide the explicit geometric witnesses, inequalities and complete proofs.

% Suggested main-text replacement/extension for the old positive-infimum theorem.
% The supplement fragment defines supp:projection-boundary.
Finite-origin positivity does not imply a globally positive lower bound. A sequence of strictly positive values may converge to zero along increasingly distant centers. This distinction is fundamental when the goal is to establish an irreducible error for \emph{every} possible normalization origin. Viewed from afar, rays through a bounded population become almost parallel, so their retained differences are associated with projections perpendicular to the viewing direction. The statistical statement needed here concerns the optimal classification risks of those representations rather than a visual resemblance between rays. The uniform theorem below controls all viewing directions simultaneously and completes the origin-search problem at infinity.

\paragraph{Uniform projection limit at distant origins.}
Fix a reference point \(p\), a unit viewing vector \(e\), and the orthogonal projection \(P_e\) onto the hyperplane perpendicular to \(e\). For binary populations with densities in dimension \(n\geq2\), the optimal risk of direction normalization about \(p+Re\) approaches the risk of \(P_e\) uniformly over \(e\) as \(R\) increases:
\begin{equation}
 \lim_{R\to\infty}\sup_e
 \left|R^*_{T_{p+Re}}-R^*_{P_e}\right|=0.
 \label{eq:projection-limit-main}
\end{equation}
Write $\Delta_\infty(e)=R^*_{P_e}-R_X^*$. The loss surface, completed by these projection limits, is continuous on a compact domain. Consequently,
\begin{equation}
 \inf_c\Delta(c)>0
 \quad\Longleftrightarrow\quad
 \begin{cases}
 \Delta(c)>0&\text{all finite }c,\\
 \Delta_\infty(e)>0&\text{all directions }e.
 \end{cases}
 \label{eq:global-criterion-main}
\end{equation}
Thus optimizing the origin cannot make loss vanish when both finite radial views and all limiting linear projections lose decisions.

\emph{Half-moon consequence.}
The connected winner regions in Theorem~1 have overlapping projection intervals in every direction. Equations~\eqref{eq:projection-limit-main}--\eqref{eq:global-criterion-main} therefore strengthen its conclusion to
\begin{equation}
 \inf_{c\in\mathbb R^2}\Delta(c)>0
 \quad\text{for }0<h\leq\sqrt{3/5}\text{ and for }h=1.
 \label{eq:global-positive}
\end{equation}
This holds for every fixed pair of positive priors and densities positive almost everywhere on their respective supports. The proof controls signed probability mass in total variation, so it establishes convergence of the optimal decision risk. Supplementary Section~\ref{supp:projection-boundary} proves the general result and gives an explicit convergence bound for bounded-support densities of bounded variation.

% Standalone main-text replacement paragraph. Definitions of d and Delta precede it.
\paragraph{A numerical lower certificate for the whole plane.}
A certificate converts the qualitative positivity theorem into a numerical lower bound valid for every origin. Such a bound can be constructed from two interior regions where one class is favored by a fixed probability margin. Their relative positions make conflicting rays unavoidable, while their sizes bound the amount of probability merged on those rays. In the plane, let $K_+,K_-$ be compact connected cores whose radius-$\rho$ neighborhoods satisfy $d\geq\kappa_+>0$ and $d\leq-\kappa_-<0$, respectively. Suppose each core meets the other's convex hull, their convex hulls contain a common disk $B(p,\delta)$, and both cores lie within distance $M$ of $p$. Writing $\kappa=\min(\kappa_+,\kappa_-)$ gives
\begin{equation}
 \inf_c\Delta(c)\geq\kappa\min\left\{
 \frac{\sqrt3\rho^3}{6M},\;2\rho\delta\right\}>0.
 \label{eq:constructive-global}
\end{equation}
For bounded origins, mutual convex-hull intrusion forces rays through both cores; their neighborhoods provide an angular interval with conflicting class mass. For distant origins, the shared convex-hull disk supplies the bound. Thus the certificate does not rest on a finite search grid. Supplementary Section~\ref{supp:constructive-global} proves the result and optimizes the split between these regions. For equal-prior area sampling, the conservative whole-plane certificates exceed $0.0205$ percentage points (pp) at $h=0.5$ and $0.00119$~pp at $h=1$. They bound $\inf_c\Delta(c)$, the best loss over \emph{all} origins; by comparison, the midpoint loss at $h=0.5$ is $10.7479\%$ for \emph{one fixed} origin (Table~\ref{tab:neural-main}). Eroded half-annuli supply the required cores under both sampling laws and any fixed positive priors.

\paragraph{Computing the best possible error at a chosen origin.}
An ideal directional classifier needs no neural training: for each viewing angle, one can integrate the class probabilities carried by that direction and choose the more probable label. For the half-annuli, circle tangencies and boundary intersections partition directions into finitely many cells. Radial support order is fixed inside a cell, so one representative ray determines whether opposite preferences coexist. Exact radial primitives then reduce Bayes risk to one angular integral. Ball arithmetic and a bounded-variation truncation bound certify its value. Supplementary Sections~\ref{supp:finite}--\ref{supp:bv} give the finite preservation test, integration formulas and error bounds. At equal priors, the density ratio $f_i^R/f_i^A\in[5/6,5/4]$ also gives $\Delta_R(c)\geq(5/6)\Delta_A(c)$.

\paragraph{Recovering the decision with one radius bit.}
The all-origin loss result rules out fixing the midpoint case by moving the origin, but it leaves open whether a small amount of distance information can repair the representation. The proposed bit answers only whether an observation is near or far from the fixed midpoint; it does not report the full distance. At $h=0.5$, direction alone has positive loss at every origin. At $c=(1/2,1/4)$, append $q=\mathbf1\{\|x-c\|>4/5\}$. For $q=0$, the sign of $9u_1+10u_2$ determines the class; for $q=1$, the sign of $u_2-u_1$ determines it. Positive signs identify class~0. Figure~\ref{fig:bit-repair} shows how the radial bit assigns observations to the two angular rules. Together, the two rules correctly label the complete class supports under their respective bit conditions, giving zero Bayes risk under either sampling law and any positive priors. Supplementary Section~\ref{supp:one-bit} proves the separating inequalities. Since zero additional radial bits leave the positive direction-only floor, one bit is the minimal radial addition at this geometry.

\paragraph{How accurately must radius be measured?}
With the same angular rules, recovery is exact precisely for $\alpha\leq t\leq\beta$, where $\alpha\approx0.7433034374$ and $\beta\approx0.8120914734$ (Threshold Addendum, Proposition~1). If the measured radius is $\widehat r=r+e$, $|e|\leq\varepsilon$, every support label is preserved for every admissible error function exactly when
\begin{equation}
\alpha+\varepsilon\leq t\leq\beta-\varepsilon.
\label{eq:radius-error}
\end{equation}
At the experimentally fixed $t=0.8$, the decoder tolerates bounded radius error up to $\beta-0.8\approx0.0120914734$ coordinate units. The different threshold $t_\star=(\alpha+\beta)/2\approx0.7776974554$ maximizes that guarantee at $(\beta-\alpha)/2\approx0.0343940180$ units (Addendum, Corollary~1). The neural experiment uses exact-radius bits at $t=0.8$; it does not test noisy radius measurements.


\begin{figure}[htbp]\centering
\includegraphics[width=.90\textwidth]{Fig2_One_Bit_Recovery.pdf}
\caption{Exact one-bit recovery for the disjoint \(h=0.5\) configuration. A shared near/far radius bit about \(c=(0.5,0.25)\) chooses between two angular decision rules. Inside threshold \(0.8\), the rule uses the sign of \(9u_1+10u_2\); outside it uses the sign of \(u_2-u_1\). The same cutoff applies to every direction. The decoder classifies the entire continuous support correctly, while the numerical neural experiment evaluates whether a fitted model learns to use this representation.}
\label{fig:bit-repair}
\end{figure}

\section{Experimental methodology and reproducibility}
\subsection{Separating a representation limit from neural learning}
The theoretical Bayes risk answers what the best classifier could do \emph{if} the class probabilities were known and the representation were fixed. A neural experiment answers how close a finite-data learner comes to that risk. These are separate measurements: an error forced by discarded radius remains even with unlimited model training, whereas error above the Bayes floor may reflect insufficient training data, architecture approximation or optimization. The empirical design therefore calculates or certifies the population risk independently and compares every trained network with the floor appropriate to its own inputs.

To support paired comparisons, the same sampled observations are transformed into each representation within a given geometry and seed. This reduces variation caused merely by generating different test sets. The study uses ten final dataset/initialization seeds, numbered 1001--1010, for the main experiment. For each seed it generates balanced training, validation and held-out test sets of 2,000, 2,000 and 50,000 observations. Distinct random streams separate dataset construction, train/validation/test splits and geometries. The dataset/initialization seed is the inferential unit; individual test points within one seed are not treated as independent experimental replicates for confidence intervals.

\subsection{Feature construction and controlled training}
The five input conditions in the main study are raw coordinates \(R\); direction-only features \(D_0,D_M,D_A\); and direction with full radius \(P_A\). Direction is calculated by dividing the centered vector by its norm, with the explicitly specified zero-residual convention. After transformation, every feature is standardized using only its training-set mean and population standard deviation; a constant feature receives unit scale. Because this transformation is invertible on each nonconstant coordinate, it does not change the representation's information content.

The principal learner is a neural network with one hidden layer of 64 hyperbolic-tangent units and a linear output. Training targets are labels \(-1\) and \(+1\), and the fitted loss is mean squared error. Full-batch Adam \citep{Kingma2015} uses \((\beta_1,\beta_2)=(0.9,0.999)\), numerical \(\epsilon=10^{-8}\), and learning rates \(0.001\) and \(0.01\). Each candidate is run for a fixed budget of 4,000 optimizer updates. Validation is performed at initialization and every 50 updates, and the selected checkpoint minimizes validation misclassification error, breaking ties by validation mean squared error, earlier step and lower learning rate.

The hidden-layer width and training protocol are the same across the main conditions; the number of parameters is not identical when feature dimension changes. A two-input model has 257 parameters, and a three-input model has 321. The population-risk comparisons isolate the loss from the representation itself. The empirical comparison of fitted networks retains this parameter-count difference as part of its stated model specification.

\subsection{Candidate selection, test isolation and diagnostic arms}
The principal factorial design specifies 300 candidates. At \(h=1\), prespecified width-15 and unconditioned width-64 diagnostics add 60 candidate fits each for \(R,D_A,P_A\). The main final experiment therefore comprises 420 candidate fits and 210 selected models. These final fits are distinct from a preceding 48-fit development stage using seeds 101 and 103. Protocols were fixed before generating and evaluating the final test sets, so their reported errors are held-out outcomes rather than checkpoint-selection statistics.

The one-bit repair experiment later adds 80 new candidate fits and 40 selected evaluations using independent random namespaces and ten fresh seeds. Accordingly, the final analytical/empirical report covers 500 candidate fits and 250 selected evaluations across the original and repair experiments. The main study and the repair study are separate evidence streams; their counts are not pooled into a fictitious single factorial comparison.

\subsection{Outcomes and statistical comparisons}
The primary outcome is test misclassification probability. Accuracy and mean squared error are secondary, while learned excess is the fitted test error minus the representation-specific Bayes risk. Six planned origin contrasts compare \(D_M-D_0\) and \(D_A-D_0\) at each of the three displacements. Paired seed differences are summarized by their means, sample standard deviations and Student-\(t\) intervals with nine degrees of freedom. Simultaneous 95\% intervals use a Bonferroni adjustment across the six contrasts; supplementary analyses retain whole-seed bootstrap sensitivity and pointwise intervals for diagnostics.

The analytical Bayes rule is also evaluated on the identical held-out observations, allowing test-sampling variation to be distinguished from the numerical integration used to calculate the population floor. Across the 15 main conditions, the mean empirical Bayes-rule error differs from the integrated risk by at most \(0.061\) percentage points. This is a consistency check of the population calculations and sampled observations, not evidence that any fitted network attained the true Bayes rule exactly.
